\section{Capítulo~\ref{sec:synthesis}. A máquina inteira}

O capítulo de síntese não traz experimentos novos: cada linha dele se apoia no capítulo em que a afirmação foi analisada e nas mesmas fontes; o barramento \nt{GLUT} e o controle de fluxo por \nt{GABA} estão firmemente estabelecidos, enquanto os papéis dos canais de difusão e o seu funcionamento conjunto são, em grande parte, hipótese.

{\footnotesize
\setlength{\tabcolsep}{3pt}\renewcommand{\arraystretch}{1.15}
\begin{longtable}{>{\raggedright}p{0.5cm}>{\raggedright}p{4.0cm}>{\raggedright}p{5.6cm}>{\raggedright}p{2.3cm}>{\raggedright\arraybackslash}p{1.7cm}}
\toprule
N.º & Afirmação & Experimento e o que foi medido & Fonte & Status \\
\midrule\endfirsthead
\toprule
N.º & Afirmação & Experimento e o que foi medido & Fonte & Status \\
\midrule\endhead
1 & Para $8.6\cdot10^{10}$ neurônios há apenas quatro tipos de sinais de controle~\eqref{eq:machine} & Contagem de núcleos celulares num homogeneizado de cérebro (fracionador isotrópico): um homem adulto tem $86.1\pm8.1$ bilhões de neurônios & \cite{Azevedo2009} & medido \\
2 & Proposição~\ref{prop:architecture}, as duas primeiras camadas: os dados correm pelo barramento de endereços \nt{GLUT}, o sinal da conexão é definido pelo emissor, o fluxo é dirigido e normalizado pelos neurônios \nt{GABA} (capítulos~\ref{sec:neurontypes}, \ref{sec:gaba}) & Um neurotransmissor principal por neurônio (revisão de medições); a inibição direta estreita a janela em que o neurônio piramidal soma as entradas; a divisão pela atividade total foi medida em muitas regiões & \cite{Strata1999,Pouille2001,Carandini2012} & medido \\
3 & Os elementos não são confiáveis: a sinapse perde a maior parte dos disparos (tab.~\ref{tab:computer}, capítulo~\ref{sec:code}) & Fatias de hipocampo, estimulação mínima: mais da metade dos disparos não provoca resposta em CA1; falhas de limiar e de condução no axônio foram excluídas, a causa é a liberação probabilística do neurotransmissor & \cite{AllenStevens1994} & medido \\
4 & O cérebro funciona com $\sim20$ watts, em média cerca de um disparo por segundo por neurônio (capítulo~\ref{sec:code}); os engenheiros ainda não construíram uma máquina assim & Estimativa~\eqref{eq:energy} a partir dos custos medidos: cerca de $10^9$ moléculas de ATP por disparo, incluindo o trabalho das sinapses (córtex de roedores); uma estimativa detalhada para o córtex dá uma taxa ainda menor. O estado da técnica segundo uma revisão & \cite{Attwell2001,Lennie2003,MehonicKenyon2022} & teoria \\
5 & O aprendizado da maioria das sinapses é o produto de um traço local por um único número comum $\delta$ (segunda equação de~\eqref{eq:machine}, capítulo~\ref{sec:dofa}) & Os neurônios \nt{DOPA} do macaco respondem ao erro de predição de recompensa; em fatias de estriado, a dopamina aumenta a espinha só se chegar $0.3$--$2$\,s depois da entrada glutamatérgica: o traço foi medido & \cite{Schultz1997,Yagishita2014} & medido \\
6 & Um fio de erro próprio para cada saída só existe no circuito de aprendizado motor (capítulo~\ref{sec:motorlearning}) & Cerebelo: a coincidência do sinal da fibra trepadeira com uma entrada enfraquece essa entrada; no macaco, os disparos complexos das células de Purkinje carregam a direção do erro de movimento e provocam a correção & \cite{Ito2001,Herzfeld2018} & medido \\
7 & Cada canal de difusão é um feixe de algumas linhas (proposição~\ref{prop:bundle}) & Para \nt{DOPA}: na mesma tarefa, neurônios diferentes carregam recompensa, movimento e atributos do estímulo; o erro rápido e o nível lento de dopamina são distintos. Para \nt{SERO}, \nt{NORA}, \nt{ACET} essa divisão foi menos estudada & \cite{Engelhard2019,Hamid2016,Mohebi2019} & medido \\
8 & \nt{SERO}: regulador de nível e, por hipótese, o horizonte $\tau_\gamma$ (capítulo~\ref{sec:serocomputer}) & Autoinibição: agonistas dos próprios receptores 5-HT$_{1A}$ inibem os neurônios serotoninérgicos da rafe (medido). O horizonte foi proposto em teoria; o experimento mais forte: a ativação optogenética desses neurônios no camundongo prolonga a espera por uma recompensa adiada & \cite{Sprouse1987,Doya2002,Miyazaki2014} & hipótese \\
9 & \nt{NORA}: o ganho $g$ escolhe entre explorar e decidir (capítulo~\ref{sec:nora}) & Aumento da relação sinal-ruído: no macaco acordado, a noradrenalina no córtex auditivo suprime a atividade de fundo mais do que a resposta a vocalizações de coespecíficos (medido); o ganho multiplicativo é um modelo; o papel na escolha entre explorar e decidir é teoria & \cite{Foote1975NE,ServanSchreiber1990,AstonJones2005} & hipótese \\
10 & \nt{ACET}: alterna entre escrita e leitura (capítulo~\ref{sec:acet}) & Três ações (enfraquecimento das conexões internas, reforço da entrada, facilitação da plasticidade) foram medidas; a troca de modos tem apoio indireto num experimento com escopolamina em humanos & \cite{HasselmoBower1992,Gil1997,Atri2004} & hipótese \\
11 & A fórmula~\eqref{eq:machine} descreve a máquina inteira & É um resumo dos modelos dos capítulos~\ref{sec:glutcomputer}--\ref{sec:acet}, cada parte se apoia no seu capítulo; como um todo, não foi testada experimentalmente & dedução no capítulo & hipótese \\
12 & Os botões trabalham de forma coordenada sem controle comum: erro, surpresa, escrita, retorno (fig.~\ref{fig:timeline}) & Não há experimento: o esquema é qualitativo. Elos isolados: a teoria da divisão de trabalho entre \nt{NORA} e \nt{ACET} e a relação entre pupila e taxa de aprendizado em humanos & \cite{YuDayan2005,Nassar2012} & hipótese \\
13 & O aprendizado por um único sinal comum é tanto mais lento quanto mais neurônios influem no resultado (proposição~\ref{prop:broadcastcost}) & Cálculo das curvas de aprendizado de uma rede linear: o aprendizado por perturbação das saídas é mais lento que a descida de gradiente tantas vezes quantas são as saídas & \cite{Werfel2005} & teoria \\
14 & Não se sabe como o córtex ajusta circuitos de muitas camadas; candidatos: rajadas de disparos, cálculos nos ramos do neurônio, autoaprendizado por predição & Rajadas como portadoras do erro: modelo; reproduzir a resposta de um modelo detalhado de neurônio cortical só uma rede de $5$--$8$ camadas conseguiu: modelo; disparos de cálcio nos ramos de neurônios do córtex humano, que dão o OU exclusivo: medido em fatias; autoaprendizado: revisão de evidências & \cite{Payeur2021,Beniaguev2021,Gidon2020,Keller2018} & hipótese \\
\bottomrule
\end{longtable}}
